\section{Aggregate Balanced Growth}

\begin{frame}{From Structural Change to Aggregate Growth}
  \small
  \begin{block}{Additional production restriction}
    To study a steady aggregate path, the paper specializes the common
    production function to Cobb-Douglas:
    \begin{equation*}
      F(n_i k_i,n_i)=k_i^\alpha n_i,\qquad \alpha\in(0,1).
    \end{equation*}
  \end{block}

  \vspace{0.15cm}
  \begin{itemize}
    \item Hicks-neutral technology with balanced aggregate growth requires a
      Cobb-Douglas form.
    \item The aggregate economy is measured in manufacturing-good units.
    \item Sectoral TFP growth enters the Euler equation through $\gammabar$.
  \end{itemize}
\end{frame}

\begin{frame}{Proposition 3: Aggregate Dynamics}
  \small
  Given any initial $k(0)$, equilibrium aggregate dynamics satisfy:
  \begin{block}{Capital accumulation}
    \begin{equation*}
      \frac{\dot k}{k}
      =A_m k^{\alpha-1}-\frac{c}{k}-(\delta+\nu).
    \end{equation*}
  \end{block}

  \begin{alertblock}{Euler equation}
    \begin{equation*}
      \theta\frac{\dot c}{c}
      =
      (\theta-1)(\gamma_m-\gammabar)
      +\alpha A_m k^{\alpha-1}-(\delta+\rho+\nu).
    \end{equation*}
  \end{alertblock}

  \vspace{0.10cm}
  The new term is $(\theta-1)(\gamma_m-\gammabar)$: it is zero in a one-sector
  model and generally time-varying under structural change.
\end{frame}

\begin{frame}[label=balancedgrowth]{Proposition 4: Conditions for Balanced Growth with Structural Change}
  \small
  \begin{alertblock}{Necessary and sufficient conditions}
    An aggregate balanced growth path with structural change exists if and only if
    \begin{equation*}
      \theta=1,\qquad \varepsilon\neq 1,\qquad
      \exists i\in\{1,\ldots,m-1\}: \gamma_i\neq\gamma_m .
    \end{equation*}
  \end{alertblock}

  \vspace{0.12cm}
  \begin{columns}[T,onlytextwidth]
    \begin{column}{0.48\textwidth}
      \begin{block}{Why $\varepsilon\neq1$?}
        Without nonunit substitution across goods, relative price changes do
        not generate employment-share movements on a balanced path.
      \end{block}
    \end{column}
    \begin{column}{0.48\textwidth}
      \begin{block}{Why $\theta=1$?}
        Log utility makes $\psi=0$, so the aggregate Euler equation no longer
        inherits the time-varying $\gammabar$ term.
      \end{block}
    \end{column}
  \end{columns}
  \hyperlink{appendixbalanced}{\beamerbutton{Intuition}}
\end{frame}

\begin{frame}{Proposition 5: Long-Run Employment Shares}
  \small
  Let sector $l$ be the limiting sector:
  the smallest TFP-growth sector if $\varepsilon<1$, or the largest TFP-growth
  sector if $\varepsilon>1$.

  \vspace{0.12cm}
  \begin{alertblock}{Asymptotic allocation on the balanced path}
    \begin{equation*}
      n_m^*=\hat{\sigma}
      =
      \alpha\left(
        \frac{\delta+\nu+g_m}{\delta+\nu+\rho+g_m}
      \right),
      \qquad
      n_l^*=1-\hat{\sigma}.
    \end{equation*}
  \end{alertblock}

  \begin{itemize}
    \item Manufacturing retains labor because it produces capital goods.
    \item The limiting consumption sector absorbs all consumption-goods labor.
    \item Other consumption sectors may decline monotonically or display
      hump-shaped employment shares.
  \end{itemize}
\end{frame}
